\section{Vulnerability management and discovery (the HW1 domain)}
\label{sec:discovery}

The planning agent of the first three sections (\cref{sec:machines,sec:records,sec:routing})
receives one scanner alert and patches it without asking whether the alert points to a flaw that an
attacker could use. The first homework, HW1, asks an agent to answer that question for every alert,
and designing such an agent requires a precise statement of the task. This section places the task
within vulnerability management, describes how discovery tools produce alerts, and introduces what
is needed to judge an alert: the attack surface, the threat model and the catalogues that classify
flaws. \Cref{sec:triage} then turns to Bandit's alerts and to triage itself.

\subsection{Vulnerability management and its stages}
\label{sec:discovery-stages}

\textbf{Vulnerability management} is the work of finding the flaws in a program that an attacker
could use, deciding which of them are real, proving that they are real, and fixing them. We use the
plain word \emph{flaw} for a defect of this kind. \Cref{sec:triage} defines the more precise term
\emph{vulnerability}.

The work has four stages, and each stage takes the result of the one before it:
\bi
\item \textbf{Discovery.} Discovery turns a repository and a scope, the part of the code to be
examined, into alerts. As in Lecture~2, an alert is a report that names a file, a line and the kind
of flaw suspected there.
\item \textbf{Triage.} Triage turns each alert into a label, together with the evidence for it. A
\textbf{label} is the verdict on one alert. In HW1 it is 1 when an attacker, as the threat model
defines one (\cref{sec:discovery-threat}), can reach the flaw the alert points to, and 0 when no
such attacker can.
\item \textbf{Validation.} Validation proves that a flaw kept by triage is real by reproducing it,
for example by running the program on an input that makes the flaw occur.
\item \textbf{Patching.} Patching changes the code so that a proven flaw is gone and the program
still behaves as it did for legitimate use.
\ei

The stages narrow the work step by step. Discovery may report many alerts, triage keeps the ones an
attacker can reach, validation confirms them, and patching removes them. A mistake in one stage
carries into the next. An alert that triage wrongly keeps costs work later, and a real flaw that
triage wrongly drops is never patched.

HW1 builds the first two stages. Its discovery stage is a fixed scan with one tool
(\cref{sec:discovery-alerts}), and the agent you build in HW1 performs triage. Validation and
patching come later in the course, with HW2 and HW3.

The planning agent works at the last of these stages, patching. HW1 adds the triage stage that the
planning agent skips.

Validation in this sense names a stage of vulnerability management. It is distinct from the planning
agent's validator (Lecture~2, Section~4), which judges whether a submitted patch should be accepted.

\subsection{Kinds of discovery tool}
\label{sec:discovery-tools}

Tools discover flaws in different ways, and each way misses some flaws. A discovery tool reads the
program's code, runs the program and watches how it behaves, or checks the libraries that the
program depends on. Static application security testing (\textbf{SAST}) reads the code. Dynamic
application security testing (\textbf{DAST}) and \textbf{fuzzing} run the program, and
\textbf{dependency scanning} checks its libraries. \Cref{tab:discovery-tools} summarizes how each
kind finds flaws, and the list after it gives what each kind misses.

\begin{table}[htbp]
  \caption{Kinds of discovery tool and how each one finds flaws.}
  \label{tab:discovery-tools}
  \begin{adjustbox}{width=\linewidth}\begin{tabularx}{\linewidth}{L{0.3\linewidth}X}
    \toprule
    \textbf{Kind} & \textbf{How it finds flaws} \\
    \midrule
    \textbf{SAST} & Matches rules against the source code, without running the program \\
    \textbf{DAST} & Runs the program and attacks it from outside, watching how it responds \\
    \textbf{Fuzzing} & Feeds the program generated inputs and keeps the inputs that make it crash \\
    \textbf{Dependency scanning} & Checks the program's libraries against published records of flaws \\
    \bottomrule
  \end{tabularx}\end{adjustbox}
\end{table}

Each kind sees some flaws and misses others:
\bi
\item \textbf{SAST.} A SAST tool reads the code without running it. Lecture~2 called a tool of this
kind a static-analysis scanner. It can examine every line, including code that no test reaches. It
never sees the values that arrive while the program runs, so where a value comes from must be
inferred from the code. Bandit's rules make no such inference, because each one matches only the
shape of the code (\cref{sec:triage}).
\item \textbf{DAST.} A DAST tool runs the program, sends it requests from outside, as an attacker
would, and watches its responses. It finds only flaws in the code that its requests cause to run.
\item \textbf{Fuzzing.} Fuzzing is dynamic testing in which the inputs are generated
automatically instead of crafted by hand. The tool keeps the inputs that make the program crash.
Like DAST, it finds only flaws in the code that its inputs cause to run.
\item \textbf{Dependency scanning.} Dependency scanning, also called software composition analysis,
compares the libraries a program uses, and their versions, with published records of known flaws.
It finds only flaws that are already known in those libraries. Flaws in the program's own code are
outside its view.
\ei

HW1 discovers flaws with one SAST tool, Bandit, a static-analysis scanner for Python, which also
produced the worked alert of Lecture~2 \citep{bandit}. HW2 and HW3 add other kinds of discovery
tool.

\subsection{Rules and what they match}
\label{sec:discovery-rules}

A SAST tool finds flaws with rules. A \textbf{rule} consists of a pattern, which is the shape of
code to look for, a message to print when the pattern matches, and metadata that classifies the
match.

% Defined here, ahead of the text that discusses it, so that this two-column float reaches the top
% of the next page, where that text runs.
\begin{figure*}[t]
  \centering
  \begin{adjustbox}{max width=\textwidth}
  \begin{tikzpicture}[
    card/.style={draw=ibmfg!60, rounded corners=2pt, fill=white, inner sep=5pt, align=left,
                 font=\sffamily\small},
    note/.style={font=\sffamily\footnotesize, text=black!60, align=center},
    bind/.style={-{Stealth[length=2mm]}, thick, draw=ibmblue},
    emit/.style={-{Stealth[length=2mm]}, thick, draw=ibmfg!80, rounded corners=3pt},
  ]
    \node[card, text width=58mm] (rule) at (0,0) {%
      \textbf{bandit.B608}, the port of Bandit's rule\\[3pt]
      \texttt{pattern-inside:}\\
      \texttt{\ \ \textcolor{ibmblue}{\$SQL} = f"...\{\textcolor{ibmblue}{\$X}\}..."}\\
      \texttt{pattern:}\\
      \texttt{\ \ \textcolor{ibmblue}{\$DB}.execute(\textcolor{ibmblue}{\$SQL})}\\[3pt]
      cwe: CWE-89\quad owasp: A1: Injection};
    \node[note, above=1mm of rule] {the rule: a pattern with gaps};
    \node[card, text width=88mm, anchor=north west] (code) at ($(rule.north east)+(24mm,0)$) {%
      \textbf{BenchmarkTest00099}\\[3pt]
      \texttt{\textcolor{black!50}{46} \textcolor{ibmblue}{sql} = f\textquotesingle SELECT username from USERS}\\
      \texttt{\phantom{46} where password = \textbackslash\textquotesingle\{\textcolor{ibmblue}{bar}\}\textbackslash\textquotesingle\textquotesingle}\\
      \texttt{\textcolor{black!50}{49} \textcolor{ibmblue}{cur}.execute(\textcolor{ibmblue}{sql})}};
    \node[note, above=1mm of code] {the code that fills the gaps};
    \node[note, below=1mm of code] (binding)
      {\texttt{\$SQL} = \texttt{sql},\quad \texttt{\$X} = \texttt{bar},\quad \texttt{\$DB} = \texttt{cur}};
    \node[card, text width=78mm, anchor=north] (alert) at ($(rule.south)+(18mm,-15mm)$) {%
      \textbf{alert}\quad bandit.B608 at line 49\\
      CWE-89 SQL injection, A1: Injection, severity WARNING};
    \node[note, below=1mm of alert] {the alert, with its kind and family of flaw attached};
    \draw[bind] (rule.east |- code.west) -- node[note, above] {binds} (code.west);
    \draw[emit] (binding.south) |- node[note, pos=0.3, right] {matches} (alert.east);
  \end{tikzpicture}
  \end{adjustbox}
  \caption{A Semgrep rule matching a BenchmarkPython test case. The metavariables of GitLab's port of
  Bandit's rule B608 bind to the variables of lines 46 and 49, and the match is reported as an alert
  that carries the rule's CWE number and Top 10 family.}
  \label{fig:discovery-binding}
\end{figure*}

Bandit's own rules are part of its Python source code (\cref{sec:triage}). Semgrep is another SAST
tool, whose rules are written as short text files, so the parts of a rule can be read directly
\citep{semgrep-rule-syntax}. GitLab has rewritten Bandit's rules in Semgrep's format
\citep{gitlab-semgrep-bandit}. A rule rewritten for another tool is called a port.

A Semgrep pattern is a piece of code with gaps in it. A gap is written as a \textbf{metavariable}, a
name that begins with \texttt{\$} and contains only capital letters, digits and underscores, such as
\texttt{\$X}. A metavariable matches one element of code in its position, such as a variable, an
expression or the name of a function. Three dots, called an ellipsis, match a sequence of zero or
more items: the arguments of a call, statements, or the characters inside a string.

When a pattern matches, each metavariable is bound to the code that fills its gap, and the tool
reports an alert at the matched line. Semgrep calls an alert a finding.

\Cref{lst:discovery-b610} shows the port of Bandit's rule B610, with two of its three patterns and
the end of its message left out.

\begin{lstlisting}[basicstyle=\ttfamily\footnotesize, caption={Bandit's rule B610 as GitLab ported it to Semgrep, shortened.}, label=lst:discovery-b610]
- id: bandit.B610
  patterns:
  - pattern-either:
    - pattern: $X.objects.$FUNC(...).extra(...)
    # two longer patterns omitted
  message: This is a last resort. You should be careful when using QuerySet.extra due to SQLi ...
  metadata:
    cwe: "CWE-89: Improper Neutralization of Special Elements used in an SQL Command ('SQL Injection')"
    owasp: 'A1: Injection'
  severity: ERROR
  languages: [python]
\end{lstlisting}

The rule has the following parts:
\bi
\item \textbf{Identifier.} The field \texttt{id} names the rule after the Bandit rule it ports.
\item \textbf{Pattern.} The block \texttt{patterns} gives the shape of code to match, and
\texttt{pattern-either} means that any one of the listed patterns may match. This pattern matches a
call to \texttt{extra} at the end of a database query written with Django, a Python web framework.
In it, \texttt{\$X} and \texttt{\$FUNC} stand for any class and any method, and the dots stand for
any arguments. The method \texttt{extra} lets a query include raw SQL text.
\item \textbf{Message.} The field \texttt{message} is the text the tool prints with each match. Here
it warns that \texttt{extra} can lead to SQL injection.
\item \textbf{Metadata.} The field \texttt{cwe} gives the kind of flaw by its number in the Common
Weakness Enumeration, a public catalogue of kinds of flaw. In this rule it is CWE-89, SQL
injection. The field \texttt{owasp} names the flaw's family: the entry of the OWASP Top 10, a second
public list, to which the flaw belongs. \Cref{sec:discovery-top10,sec:discovery-cwe} describe both
lists.
\item \textbf{Severity and language.} The field \texttt{severity} is the rule's estimate of the harm
if a match is a real flaw. Semgrep grades it with words such as ERROR and WARNING, and Bandit grades
it as LOW, MEDIUM or HIGH (\cref{sec:triage}). The field \texttt{languages} names the programming
languages whose files the rule applies to.
\ei
A rule therefore carries the shape to match, the message to print, and the kind and family of flaw
that a match belongs to.

\Cref{fig:discovery-binding} shows a rule binding to real code. The rule is GitLab's port of
Bandit's B608, which looks for SQL queries built from strings. Its \texttt{pattern-inside} clause
requires that the match lie inside code where a variable \texttt{\$SQL} is assigned an f-string with
a gap \texttt{\{\$X\}}. Inside the string, the dots match any characters, so the clause accepts any
f-string with a gap. Its \texttt{pattern} clause then requires a call \texttt{\$DB.execute(\$SQL)}
with that variable as its only argument.

The code is the test case BenchmarkTest00099 of OWASP BenchmarkPython, the Python web application
that the planning agent works on (\cref{sec:machines}). As in Lecture~2, a test case is a short program
that handles one kind of web request. Line~46 builds a query as an f-string with \texttt{bar} in
its gap, and line~49 executes it. The names set in blue are the metavariables and the code they
bind to: \texttt{\$SQL}~=~\texttt{sql}, \texttt{\$X}~=~\texttt{bar} and \texttt{\$DB}~=~\texttt{cur}.
The port reports an alert at line~49 with CWE-89, the family tag A1: Injection, and the severity
this rule assigns, WARNING.

A rule can hold several patterns, one for each way of writing the same flaw. Besides the form of
\cref{fig:discovery-binding}, the port's B608 entry matches a query built with the \texttt{\%}
operator, with \texttt{.format}, with an f-string or by concatenation directly inside the call to
\texttt{execute}. It also matches a query built earlier with \texttt{\%} or \texttt{.format} and
stored in a variable.

Which code produces an alert depends only on the rule's pattern. GitLab's port writes the pattern
of B608 in Semgrep's notation, while Bandit's own B608 is Python code that checks whether the text
of a string looks like a SQL query (\cref{sec:triage}). The two patterns differ, so the two
implementations report different alerts:
\bi
\item \textbf{The reported line.} Bandit's own B608 reports the flaw of
\cref{fig:discovery-binding} at line~46, where the query is built. The port reports line~49, where
the query is executed.
\item \textbf{The matched code.} Another test case, BenchmarkTest00011, writes its query as an
f-string with no gap and passes the value to the database separately, as a bound parameter
(Lecture~2). The port does not match it. Its query has no gap for \texttt{\$X} to bind, and its
call to \texttt{execute} passes a second argument, which \texttt{\$DB.execute(\$SQL)} does not
allow. Bandit's own B608 does report it, because it flags any f-string whose text looks like a SQL
query, whether or not the f-string has a gap.
\ei

\subsection{What an alert contains}
\label{sec:discovery-alerts}

In HW1, discovery runs Bandit's fixed set of rules over the code and reports every place where a
rule matches. Each alert names a file, a line, the rule that matched and the kind of flaw that the
rule suspects, so that the next stage knows where to look. An alert records only that a rule
matched. Whether the match is a real flaw is left open, and deciding it is the job of triage.

The worked alert of Lecture~2 is an example. Bandit's rule B608 matched a line of BenchmarkPython that
builds a SQL query by interpolation, and the alert named the file, that line and the rule. The alert
did not say whether a value supplied by a user could reach the query. BenchmarkTest00099, the test
case of \cref{fig:discovery-binding}, is a different test case. It builds and runs its query with the
same two lines as the worked alert's test case, at lines~46 and~49, and Bandit reports it in the
same way, at line~46.

HW1 scans three targets with Bandit 1.9.4 and keeps the alerts that Bandit rates at MEDIUM severity
or higher. The three targets are the following:
\bi
\item \textbf{OWASP BenchmarkPython.} The first 99 test cases of BenchmarkPython give 18 alerts
\citep{owasp-benchmarkpython}.
\item \textbf{Radicale.} Radicale, a calendar server, gives 10 alerts \citep{radicale}.
\item \textbf{SGLang.} SGLang, a server that runs large language models and answers requests to
them, gives 21 alerts \citep{sglang}.
\ei

Each target is pinned at one commit, a fixed version of its code, and Bandit's rules are fixed as
well. Scanning the same commit therefore gives the same list of alerts every time. Every agent built
for HW1 receives the same alerts, and two runs of an agent can be compared alert by alert.

The targets play different roles in HW1. BenchmarkPython ships an \textbf{answer key}, a file that
records the true labels: for each test case, its category, the CWE number it tests, and whether its
flaw is real. A \textbf{development set} is a target whose answer key is open to students, so that
they can check their agent's labels while they build it. The first 99 test cases of BenchmarkPython
are HW1's development set.

The agent itself never reads the answer key (Lecture~2, Section~2). An agent that could read it would
copy its labels instead of analyzing the code. A \textbf{held-out target} is a target whose answer
key stays hidden. Radicale and SGLang are held out, so they measure how the agent performs on code
it was not tuned on.

\subsection{The attack surface}
\label{sec:discovery-surface}

An attacker can affect a program only through the data the program takes in from outside. That data
arrives at a few places, such as a request, a file or a message, and an attacker starts there. The
\textbf{attack surface} of a program is the set of places where data from outside the program enters
it, together with the code that first handles that data.

\begin{definition}[Attack surface]
\label{def:discovery-surface}
``Intuitively, a system's attack surface is the set of ways in which an adversary can enter the
system and potentially cause damage.'' \citep[\S1]{manadhata2007approach}
\end{definition}

In the formal version of this definition, a system's attack surface consists of its entry and exit
points, its channels and its untrusted data items \citep{manadhata2011attack}. For triage, the part
that matters most is where data enters. An \textbf{entry point} is one place where data from outside
enters the program.

Each HW1 target has its own attack surface:
\bi
\item \textbf{BenchmarkPython.} Every test case is a Flask route, a function that the Flask web
framework calls to handle requests to one address, and it reads a cookie or a form field from the
request. Its surface is the request.
\item \textbf{Radicale.} Radicale is a calendar server, so its surface is the HTTP request that a
client sends, whether the client has logged in or not.
\item \textbf{SGLang.} SGLang is a server for large language models. Its surface includes the
request, and also the model files it loads and the channels over which its own worker processes
exchange messages.
\ei

\Cref{lst:discovery-entry} shows one entry point in code, the start of the test case of
\cref{fig:discovery-binding}.

\begin{lstlisting}[basicstyle=\ttfamily\footnotesize, caption={The entry point of the test case BenchmarkTest00099 and the path of its value to \texttt{bar} (lines 31 to 42, with lines 34 to 37 and 40 left out).}, label=lst:discovery-entry]
31  param = request.form.get("BenchmarkTest00099")
32  if not param:
33      param = ""
38  conf52528 = configparser.ConfigParser()
39  conf52528.add_section('section52528')
41  conf52528.set('section52528', 'keyB-52528', param)
42  bar = conf52528.get('section52528', 'keyB-52528')
\end{lstlisting}

Line~31 is the entry point. It reads a form field of the incoming request, and the client who sends
the request controls its value. Lines~32 and~33 replace a missing value with an empty string.
Lines~38 to~42 store the value in a configuration object under one key and read it back under the
same key, so \texttt{bar} holds exactly what the client sent. This detour is the benchmark's way of
hiding a flow of data from a scanner.

A \textbf{sink} is the call where outside data can do harm, such as the call that sends a query to a
database. In BenchmarkTest00099, line~46 places \texttt{bar} in the query text, and line~49 sends
that text to the database (\cref{fig:discovery-binding}). Line~49 is the sink, and whatever the
client sends in the form field is the value of \texttt{bar} there. A scanner's rule usually points
at the sink or at the line that prepares its argument. GitLab's port points at line~49, and
Bandit's own B608 points at line~46 (\cref{sec:discovery-rules}).

For every alert, the agent you build in HW1 asks whether the value at the flagged line can come from
the attack surface. In this test case it can. The value read at line~31 reaches the sink at line~49
unchanged.

\subsection{The threat model}
\label{sec:discovery-threat}

A label says whether an attacker can reach a flaw, so it needs a definite attacker: who the attacker
is and which data they control. A \textbf{threat model} names the attacker, the data the attacker
controls, and the harm the attacker must be unable to do. \Cref{tab:discovery-threat} gives the
attacker of each HW1 target. Its last column, what the attacker controls, is what a label depends on.

\begin{table}[htbp]
  \caption{The attacker of each HW1 target and the data that attacker controls.}
  \label{tab:discovery-threat}
  \begin{adjustbox}{width=\linewidth}\begin{tabularx}{\linewidth}{L{0.22\linewidth}L{0.3\linewidth}X}
    \toprule
    \textbf{Target} & \textbf{The attacker} & \textbf{What they control} \\
    \midrule
    \textbf{Benchmark\-Python} & An anonymous web client & Every byte of the request: cookies, form fields, headers \\
    \textbf{Radicale} & A calendar client, logged in or not & The request, plus the storage directory if local users count \\
    \textbf{SGLang} & A client, or whoever publishes a model & The request, the model files it loads, and its worker channel \\
    \bottomrule
  \end{tabularx}\end{adjustbox}
\end{table}

% Defined in the threat-model subsection, ahead of the Top 10 text that first cites it, so that this
% two-column float reaches the top of the page on which that text runs (a table* defined on one page
% appears at the earliest on the next). It precedes the Top 10 table in the source because LaTeX keeps
% tables in number order: defined after it, it waits behind that table when the table is deferred.
\begin{table*}[t]
  \caption{From a Top 10 family to a CWE, to the Python call that carries the flaw, and to the
  BenchmarkPython category that tests it. The CWE numbers are those of the benchmark's answer key.}
  \label{tab:discovery-hierarchy}
  \begin{adjustbox}{width=\linewidth}\begin{tabularx}{\linewidth}{L{0.36\linewidth}L{0.08\linewidth}XL{0.15\linewidth}}
    \toprule
    \textbf{Family} & \textbf{CWE} & \textbf{The Python call} & \textbf{Benchmark category} \\
    \midrule
    \textbf{Injection} & CWE-89 & A query built from a string & \texttt{sqli} \\
    \textbf{Injection} & CWE-78 & A shell command built from a string & \texttt{cmdi} \\
    \textbf{Injection} & CWE-94 & \texttt{exec} or \texttt{eval} on a string & \texttt{codeinj} \\
    \textbf{Cryptographic Failures} & CWE-328 & MD5 or SHA1 applied to a secret & \texttt{hash} \\
    \textbf{Software or Data Integrity Failures} & CWE-502 & \texttt{pickle.load} or \texttt{torch.load} & \texttt{deserialization} \\
    \textbf{Security Misconfiguration} & CWE-611 & An XML parser applied to outside data & \texttt{xxe} \\
    \textbf{Broken Access Control} & CWE-22 & \texttt{open}, or a test that a file exists, on a path from outside & \texttt{pathtraver} \\
    \bottomrule
  \end{tabularx}\end{adjustbox}
\end{table*}

For BenchmarkPython the threat model is simple, because an anonymous client controls the whole
request. For Radicale and SGLang, the labels of most alerts depend on which attackers the threat
model counts:
\bi
\item \textbf{Radicale.} Seven of its ten alerts load files with Python's \texttt{pickle} module,
which turns stored bytes back into Python objects and can run code chosen by whoever wrote the file.
Radicale wrote these files itself, in its storage directory. The storage directory is in the
attacker's control only if the threat model counts local users of the machine as attackers.
\item \textbf{SGLang.} Eighteen of its 21 alerts load models, from model files or from the Hugging
Face Hub, or unpickle bytes that do not come from the HTTP request: pickled files on disk, and
messages on the channels over which SGLang's processes and their clients exchange work. Whether
they are reachable depends on whether the threat model counts as an attacker whoever publishes a
model, whoever writes those files, or whoever can send messages on those channels.
\ei

The labels of these alerts turn entirely on the threat model. Bandit ignores the threat model,
because its rules match the shape of the code and carry no information about who controls the data.

A \textbf{rationale} is the text that justifies a label. The agent you build in HW1 should state the
threat model once in its prompt, so that every label is decided under the same assumption. It should
state it again in every rationale, so that a reviewer can see which attacker the label assumed.
Without a stated threat model a label cannot be checked, because the same alert can deserve 1 under
one threat model and 0 under another.

\subsection{The OWASP Top 10}
\label{sec:discovery-top10}

The \textbf{OWASP Top 10} is a ranked list of the most critical security risks to web
applications, published by the OWASP Foundation (Open Worldwide Application Security Project). A
\textbf{family} is one of its ten entries: a group of related kinds of flaw, such as injection. The
current edition, the OWASP Top 10:2025, was announced as a release candidate on 6 November 2025 and
published in final form in December 2025 \citep{owasp2025top10}.
Eight of its families were chosen and ranked from data on over 2.8 million applications, and two
were added from a survey of practitioners. The analysis considered 589 kinds of flaw from the
catalogue described in \cref{sec:discovery-cwe}.

\Cref{tab:discovery-top10} lists the ten families with a plain description of each. The families
that matter most for HW1 are Software or Data Integrity Failures, Injection, Cryptographic Failures
and Security Misconfiguration. \Cref{tab:discovery-hierarchy} follows these families down to the
Python calls that carry their flaws, and \cref{sec:discovery-python} explains it.

\begin{table}[htbp]
  \caption{The ten families of the OWASP Top 10:2025, each with a plain description.}
  \label{tab:discovery-top10}
  \begin{adjustbox}{width=\linewidth}\begin{tabularx}{\linewidth}{>{\RaggedRight\arraybackslash\hyphenpenalty=10000\exhyphenpenalty=10000}p{0.56\linewidth}X}
    \toprule
    \textbf{Family} & \textbf{What goes wrong} \\
    \midrule
    \textbf{A01 Broken Access Control} & A user reaches data or actions that should be denied to them \\
    \textbf{A02 Security Misconfiguration} & A setting enables something that should be off \\
    \textbf{A03 Software Supply Chain Failures} & A dependency brings in unchecked code \\
    \textbf{A04 Cryptographic Failures} & Protection is missing or weak \\
    \textbf{A05 Injection} & Outside data becomes part of a command \\
    \textbf{A06 Insecure Design} & The design lacks a control it needed \\
    \textbf{A07 Authentication Failures} & The program fails to confirm who the user is \\
    \textbf{A08 Software or Data Integrity Failures} & Code or data is trusted without a check \\
    \textbf{A09 Security Logging and Alerting Failures} & An attack is not recorded, or its record is not reviewed \\
    \textbf{A10 Mishandling of Exceptional Conditions} & An error is handled in a way that grants access or reveals data \\
    \bottomrule
  \end{tabularx}\end{adjustbox}
\end{table}

The edition matters when reading a rule's metadata. The tag \texttt{A1: Injection} of
\cref{lst:discovery-b610} uses the numbering of the 2017 edition \citep{owasp2017top10}. In the 2025
edition, Injection is A05.

The 49 alerts of HW1 fall into the families as follows:
\bi
\item \textbf{Software or Data Integrity Failures.} This family has 26 alerts. They load pickled
data or model files, or download models without an integrity check.
\item \textbf{Injection.} This family has 7 alerts. They build SQL queries from strings, run
\texttt{exec} or \texttt{eval}, or start a shell command.
\item \textbf{Cryptographic Failures.} This family has 6 alerts, which use the MD5 or SHA1 hash.
\item \textbf{Security Misconfiguration.} This family has 6 alerts, which parse XML.
\Cref{sec:discovery-python} explains why XML parsing belongs to this family.
\item \textbf{None of the ten.} The remaining 4 alerts come from Bandit's rule B113, which no family
covers (\cref{sec:discovery-python}).
\ei
The integrity alerts of Radicale and SGLang are exactly the alerts whose labels turn on the threat
model (\cref{sec:discovery-threat}). The other integrity alert comes from BenchmarkPython.

\subsection{Weaknesses and their catalogues}
\label{sec:discovery-cwe}

A \textbf{weakness} is a kind of flaw that recurs across programs. Each family of the Top 10 is a
set of weaknesses from a public catalogue. The Common Weakness Enumeration (\textbf{CWE}) is that
catalogue, kept by MITRE, and it gives each weakness a number and a name \citep{mitre-cwe}. An entry
of the catalogue is also called a CWE. CWE-89, for example, is SQL injection, and it can occur in
any program, in any language, that builds a query from text.

CWE, CVE and CVSS are easily confused, because they describe flaws at different levels:
\bi
\item \textbf{CWE.} A CWE is a weakness, such as CWE-89. It describes no particular program.
\item \textbf{CVE.} A Common Vulnerabilities and Exposures record (CVE) identifies one
instance of a flaw in one product \citep{cve-program}. Its record usually names a CWE as the type
of the flaw.
\item \textbf{CVSS.} The Common Vulnerability Scoring System (CVSS) rates one instance of a
flaw by how easily an attacker can use it and how much damage it does
\citep{first2023cvss4}. A CVSS score therefore attaches to a CVE.
\ei

A family groups many CWEs. In the 2025 edition, the Injection family alone contains 37 CWEs,
including CWE-89 (SQL injection), CWE-78 (OS command injection), CWE-94 (code injection) and
CWE-95 (eval injection) \citep{owasp2025top10}.

Every HW1 alert carries a CWE number, because Bandit 1.9.4 attaches one to every alert it reports.
The number names the kind of flaw that the rule suspects at the flagged line.

\subsection{From a family to a line of Python}
\label{sec:discovery-python}

The Top 10 and the CWE catalogue form a hierarchy that continues into the code: a family contains
CWEs, and each CWE appears in code as a particular kind of Python call.
\Cref{tab:discovery-hierarchy} follows this hierarchy for several kinds of flaw that
BenchmarkPython tests. Its last column names the benchmark category whose test cases cover each
kind. The answer key has further categories, such as XPath injection (\cref{sec:triage}). Read each
row from left to right, from the general family to the specific call.

The call at the bottom of the hierarchy is the sink (\cref{sec:discovery-surface}). An alert's CWE
therefore usually tells the agent the family of the suspected flaw and the shape of code to look
for around the flagged line.

The CWE numbers in \cref{tab:discovery-hierarchy} are those of BenchmarkPython's answer key. Bandit
attaches its own numbers, which sometimes differ from the key:
\bi
\item \textbf{exec and eval.} Bandit reports CWE-78, OS command injection, where the key has CWE-94,
code injection. Both are Injection CWEs.
\item \textbf{MD5 and SHA1.} Bandit reports CWE-327, a broken or risky cryptographic algorithm, where
the key has CWE-328, a weak hash. Both belong to Cryptographic Failures.
\item \textbf{XML parsers.} Bandit reports CWE-20, improper input validation. The flaw these rules
look for is an XML external entity (XXE) flaw, in which an XML document makes the parser read a file
or contact a server that the document names. The key gives CWE-611 for its XXE cases. CWE-20 is an
Injection CWE in the 2025 edition, while CWE-611 belongs to Security Misconfiguration.
\ei

The XML alerts are the exception to this correspondence. Their CWE points to Injection, while the
flaw they look for belongs to Security Misconfiguration, which is where the counts of
\cref{sec:discovery-top10} place them. In the development set, the two XML alerts of
BenchmarkTest00017 fall in a test case of the XXE category. The XML alerts of BenchmarkTest00013 to
BenchmarkTest00016 fall in test cases of the XPath-injection category, for which the key gives
CWE-643 (\cref{sec:triage}).

A scanner's rules and the families of the Top 10 are separate lists. A rule can report a flaw that
no family contains, and a family can contain flaws that no rule reports. Bandit's rule B113 reports
an HTTP request made without a timeout, which can leave the program waiting indefinitely, and it
attaches CWE-400, uncontrolled resource consumption. CWE-400 belongs to none of the ten families,
so the four B113 alerts of HW1 have no family.

Path traversal is a flaw of the second kind. In a path traversal (CWE-22, Broken Access Control), a
request makes the program open a file outside the folder it intended. BenchmarkPython's
path-traversal cases open a file, or test whether a file exists, on a path built from the request,
and Bandit has no rule for either shape of code. The
path-traversal cases of the development set therefore receive no Bandit alert (\cref{sec:triage}).
The path-traversal alert of the second recorded run in \cref{sec:routing} came from a different
scanner.

\subsection*{Key points}
\bi
\item Vulnerability management has four stages: discovery turns code into alerts, triage turns each
alert into a label with its evidence, validation proves a kept flaw by reproducing it, and patching
removes a proven flaw while keeping the program's behavior.
\item Discovery tools read the code (SAST), run the program (DAST and fuzzing) or check its libraries
(dependency scanning), and a SAST tool such as Bandit matches rules, each made of a pattern, a
message and metadata.
\item An alert records only that a rule matched, and HW1's alerts come from one fixed Bandit scan of
a development set with an open answer key (BenchmarkPython's first 99 test cases) and of two
held-out targets (Radicale and SGLang).
\item The attack surface is the set of places where outside data enters a program, with the code
that first handles it, and triage asks whether the value at the sink can come from an entry point.
\item A threat model names the attacker, the data the attacker controls and the harm to prevent, and
because the same alert can deserve 1 under one threat model and 0 under another, the agent states it
in its prompt and in every rationale.
\item The OWASP Top 10:2025 ranks the most critical security risks to web applications as ten
families, each a set of CWEs, where a CWE is a kind of flaw, a CVE is one instance of a flaw in one
product, and a CVSS score rates that instance.
\item An alert's CWE usually names its family and the shape of code to look for, although Bandit's
XML alerts point to the wrong family and the B113 alerts belong to no family.
\ei
